% Hafta 2 — Polimorfik virüs neden imzayla yakalanmaz?
% Derleme: py -3.12 tools/latex/derle.py docs/week-2/assets/h02-04-polimorfik.tex
\documentclass[tikz,border=8pt]{standalone}
\usepackage{cen429-cizim}
\begin{document}
\begin{tikzpicture}[node distance=9mm and 14mm]
  \node[baslik] (b) at (0,0) {Polimorfik virüs neden imzayla yakalanmaz?};
  \node[altbaslik] at (b.south west) {her kopyada baytlar değişir, açılan gövde aynı kalır};

  \node[notrkutu=50mm, draw=kotu, below=20mm of b.south west, anchor=north west, xshift=10mm] (cozucu)
    {\kb{Çözücü kod}\\her kopyada FARKLI};
  \node[notrkutu=50mm, draw=kotu, below=of cozucu] (sifreli)
    {\kb{Şifreli gövde}\\her kopyada FARKLI baytlar};
  \node[iyikutu=50mm, below=of sifreli] (acik)
    {\kb{Açık gövde}\\HER ZAMAN AYNI};

  \draw[ok, draw=soluk] (cozucu) -- (sifreli);
  \draw[ok, draw=soluk] (sifreli) -- node[oketiket, right] {çöz} (acik);

  \begin{pgfonlayer}{background}
    \node[draw=kotu, fill=kotubg, rounded corners=4pt, line width=1.3pt,
          fit=(cozucu)(sifreli)(acik), inner sep=8mm,
          label={[font=\sffamily\bfseries\small, text=kotu]above:Şifreli / polimorfik virüs}] {};
  \end{pgfonlayer}

  \node[iyikutu=40mm, right=42mm of sifreli] (imza) {\kb{İmza eşleşmesi}\\burada yakalanır};
  \draw[iyiok] (acik.east) to[out=10, in=-100] node[oketiket, pos=0.55, align=center, text=iyi]
    {emülasyon\\çözünce} (imza.south);

  \node[dip, text width=170mm, below=12mm of acik.south -| cozucu, anchor=north]
    {Bu yüzden emülasyon kullanılır: virüsü yalıtılmış ortamda çalıştırıp AÇIK gövdeyi imzalar.};
\end{tikzpicture}
\end{document}
